\documentclass[border=7bp]{standalone}
% Shared preamble for the report's editable vector figures.
\usepackage{fontspec}
\usepackage{xeCJK}
\setmainfont{Calibri}
\setCJKmainfont{SimSun}[AutoFakeBold=1.5]
\usepackage{amsmath}
\usepackage{tikz}
\usetikzlibrary{arrows.meta,positioning,calc,shapes.geometric,fit}
\definecolor{Ink}{HTML}{25313D}
\definecolor{Blue}{HTML}{315E80}
\definecolor{PaleBlue}{HTML}{E7EFF5}
\definecolor{Muted}{HTML}{626B74}
\definecolor{PaleGray}{HTML}{F0F1F3}
\tikzset{
  >=Stealth,
  every node/.style={font=\fontsize{10bp}{13bp}\selectfont,text=Ink},
  flow/.style={->,draw=Ink,line width=.65pt},
  frame/.style={rectangle,draw=Ink,line width=.65pt,align=center,
    minimum height=.85cm,text width=5cm,inner sep=5pt},
  emphasis/.style={frame,draw=Blue,fill=PaleBlue},
  supporting/.style={frame,draw=Muted,fill=PaleGray},
  annotation/.style={font=\fontsize{9bp}{12bp}\selectfont,text=Muted,align=center},
}

\begin{document}
\begin{tikzpicture}

\node[emphasis,text width=11cm] (input) at (0,0) {输入：棋盘初态、目标、算法、启发函数与运行配置};
\node[frame,text width=11cm] (runner) at (0,-1.45) {运行入口：SearchTester / StudentSearchTester / 实验运行器\\教学入口由 EngineFeeder 提供问题、\mbox{Frontier} 与 Predictor；实验入口按配置直接构造};
\node[emphasis,text width=11cm] (search) at (0,-3.1) {Searcher：A* 的最小 $f$ 出队，或 IDA* 的阈值内深度优先搜索};
\node[frame,text width=4.6cm] (problem) at (-3.2,-5) {Problem / State\\目标判定、合法动作、状态转换\\生成后继并累计 $g$};
\node[frame,text width=4.6cm] (heuristic) at (3.2,-5) {Predictor：计算 $h$\\错位数 / \mbox{Manhattan} / LC / PDB\\PDB 先建表或加载，再查询};
\node[supporting,text width=11cm] (manage) at (0,-7) {候选管理与重复状态处理\\A*：Frontier + Closed；IDA*：当前路径 + 可选 TT / ZOBRIST};
\node[emphasis,text width=11cm] (solution) at (0,-8.6) {找到目标：沿 Node 父链恢复动作序列，输出解路径 TSV};
\node[frame,text width=11cm] (visual) at (0,-10.05) {SolutionTrace 读取文件 → PlaybackModel 控制步号 → PuzzleVisualizer 显示与播放};
\draw[flow] (input)--(runner); \draw[flow] (runner)--(search);
\draw[flow] (search.south)-|(problem.north); \draw[flow] (search.south)-|(heuristic.north);
\coordinate (join) at (0,-6.35);
\draw[draw=Ink,line width=.65pt] (problem.south)|-(join);
\draw[draw=Ink,line width=.65pt] (heuristic.south)|-(join);\draw[flow](join)--(manage.north);
\draw[flow] (manage.west)--++(-.85,0)|-node[annotation,pos=.23,left]{继续搜索}(search.west);
\draw[flow] (search.east)--++(1.4,0)|-node[annotation,pos=.25,right]{目标成立}(solution.east);
\draw[flow] (solution)--(visual);
\node[annotation] at (0,-11) {上半部分完成求解；下半部分通过文件衔接播放器，不在播放时重新搜索};

\end{tikzpicture}
\end{document}
